\section{High-Agreement Tangent Staircase}
\label{sec:tangent-staircase}

\textbf{Status: \textsc{New-local} finite MDS theorem.}
This section records a generic tangent floor and, in the very-high-agreement
range, a matching upper bound.  It is independent of quotient-periodic
structure: smoothness is not used.  The theorem is nevertheless directly
relevant to the prize ledger because the floor is paid with the actual line
field \(\qline\).

For a polynomial \(P\in F[X]\), write
\[
  P=P_{<k}+P_{\ge k},
  \qquad \deg P_{<k}<k,
\]
where \(P_{\ge k}\) is the sum of the terms of degree at least \(k\).  Let
\(\operatorname{LD}_{\rm sw}(C,a)\) denote the largest possible number of finite
slopes on a received line that have a support-wise noncontained explanation on
some support of size at least \(a\).

\begin{theorem}[moving-root tangent floor]
\label{thm:moving-root-tangent-floor}
Let \(D\subset F\) have \(|D|=n\), and let
\(C=\RS[F,D,k]\).  For every integer \(a\) with \(k+1\le a\le n\),
\[
  \operatorname{LD}_{\rm sw}(C,a)\ge n-a+1 .
\]
\end{theorem}

\begin{proof}
Choose a subset \(A\subset D\) of size \(a-1\), and put
\[
  B_A(X)=\prod_{x\in A}(X-x).
\]
Define the received line
\[
  f=(XB_A)_{\ge k},
  \qquad
  g=-(B_A)_{\ge k}.
\]
For every \(t\in D\setminus A\), set
\[
  L_t(X)=(X-t)B_A(X).
\]
Then \(L_t\) is the locator of the support \(A\cup\{t\}\), which has size
\(a\), and
\[
  f+t g=(L_t)_{\ge k}.
\]
On every root of \(L_t\), the word represented by \((L_t)_{\ge k}\) agrees with
\(- (L_t)_{<k}\), a codeword of \(C\).  Thus the slope \(t\) is explained on
\(A\cup\{t\}\).

It remains to check support-wise noncontainment on this same support.  On
\(A\), since \(B_A=0\), we have
\[
  g=-(B_A)_{\ge k}=(B_A)_{<k}.
\]
If \(g\) were explained on \(A\cup\{t\}\) by a degree-\(<k\) polynomial \(q\),
then \(q=(B_A)_{<k}\), because \(|A|=a-1\ge k\).  At the remaining point \(t\),
this would force
\[
  -(B_A)_{\ge k}(t)=(B_A)_{<k}(t),
\]
that is, \(B_A(t)=0\), contrary to \(t\notin A\).  Therefore the explanation is
support-wise noncontained.  The slopes \(t\in D\setminus A\) are distinct, so
there are \(n-a+1\) of them.
\end{proof}

\begin{theorem}[very-high-agreement tangent staircase]
\label{thm:very-high-tangent-staircase}
Let \(C=\RS[F,D,k]\) with \(|D|=n\).  If
\[
  3a-2n\ge k,
\]
then
\[
  \operatorname{LD}_{\rm sw}(C,a)=n-a+1 .
\]
\end{theorem}

\begin{proof}
The lower bound is Theorem~\ref{thm:moving-root-tangent-floor}.  For the upper
bound, put \(s=n-a\).  Consider any received line with at least two finite
support-wise noncontained slopes at agreement at least \(a\); the one-slope
case is trivial.  Pick two such slopes \(z_1\ne z_2\), with explaining
codewords \(p_1,p_2\in C\) on supports \(S_1,S_2\), respectively.  Since
\(|S_i|\ge a\),
\[
  |S_1\cap S_2|\ge 2a-n.
\]
On this intersection,
\[
  g=(p_1-p_2)/(z_1-z_2),
  \qquad
  f=p_1-z_1 g .
\]
Because \(2a-n\ge k\) follows from \(3a-2n\ge k\), the two right-hand sides
extend uniquely to codewords \(c_g,c_f\in C\).  Let \(S_0\) be a maximal common
support on which \(f=c_f\) and \(g=c_g\), and put \(b=|S_0|\) and
\(d=n-b\).  Then \(b\ge 2a-n\), so \(d\le 2s\).  By maximality, outside
\(S_0\) there is no coordinate at which both residuals \(f-c_f\) and
\(g-c_g\) vanish.

The residual-budget proposition, Proposition~\ref{prop:common-residual}, is
applicable because
\[
  a+b-n\ge a+(2a-n)-n=3a-2n\ge k .
\]
If \(b\ge a\), then every noncontained slope needs at least one residual zero
outside \(S_0\), and there are only \(d\le s\) such coordinates.  If \(b<a\),
then every noncontained slope needs at least
\[
  h=a-b=d-s
\]
private residual zeros outside \(S_0\).  Since \(s<d\le2s\), the residual
budget gives
\[
  \#\{\text{bad slopes}\}
  \le
  \left\lfloor \frac d{d-s}\right\rfloor
  \le s+1
  =n-a+1 .
\]
This proves the matching upper bound.
\end{proof}

\begin{corollary}[exact high-agreement gate for the \(\F_{17^{32}}\), rate-\(1/2\) row]
\label{cor:f17-32-exact-tangent-gate}
Let \(H\subset \F_{17^{32}}^*\) be any multiplicative subgroup of order
\(512\), and let
\[
  C=\RS[\F_{17^{32}},H,256].
\]
For every \(a\ge427\),
\[
  \operatorname{LD}_{\rm sw}(C,a)=513-a.
\]
In particular,
\[
  \operatorname{LD}_{\rm sw}(C,506)=7,
  \qquad
  \operatorname{LD}_{\rm sw}(C,507)=6.
\]
Since
\[
  6\cdot2^{128}<17^{32}<7\cdot2^{128},
\]
we have
\[
  \frac{\operatorname{LD}_{\rm sw}(C,507)}{17^{32}}<2^{-128}
  \quad\text{but}\quad
  \frac{\operatorname{LD}_{\rm sw}(C,506)}{17^{32}}>2^{-128}.
\]
Thus the finite-slope support-wise MCA threshold for this row is pinned at the
agreement staircase between \(506\) and \(507\).  Equivalently, the largest safe
integer Hamming radius is \(5\), while integer radius \(6\) is already unsafe.
The corresponding normalized grid radii are
\[
  5/512
  \quad\text{and}\quad
  6/512=3/256.
\]
\end{corollary}

\begin{proof}
Here \(n=512\) and \(k=256\).  The hypothesis
\(3a-2n\ge k\) becomes \(3a-1024\ge256\), hence \(a\ge427\).  The exact formula
is Theorem~\ref{thm:very-high-tangent-staircase}.  The arithmetic inequalities
are
\[
  17^{32}-6\cdot2^{128}
  =326217393234836465063858652730341978625>0
\]
 and
\[
  7\cdot2^{128}-17^{32}
  =14064973686101998399515954701426232831>0 .
\]
\end{proof}

\begin{remark}[effect on the earlier agreement-353 target]
The agreement-353 target is no longer a frontier for the rate-\(1/2\),
\(\F_{17^{32}}\), \(n=512\) row: Theorem~\ref{thm:moving-root-tangent-floor}
already gives
\[
  \operatorname{LD}_{\rm sw}(C,353)\ge512-353+1=160.
\]
The next unresolved issue is not to find seven slopes at agreement \(353\), but
to decide whether any non-tangent mechanism survives beyond the exact tangent
staircase at agreement \(507\).
\end{remark}
